\section*{Chapter \ref{ch:g10:concentration-and-dilution} --- Concentration and Dilution}

\begin{solution}{exo:g10:concentration-and-dilution:1}
$C_m = \qty{5.0}{g} / \qty{0.250}{L} = \qty{20}{g/L}$.
\end{solution}

\begin{solution}{exo:g10:concentration-and-dilution:2}
$c = \qty{0.050}{mol} / \qty{0.200}{L} = \qty{0.25}{mol/L}$.
\end{solution}

\begin{solution}{exo:g10:concentration-and-dilution:3}
$M(\ce{CuSO4}) = 63.5 + 32.1 + 4 \times 16.0 = \qty{159.6}{g/mol}$;
$n = 0.10 \times 0.1000 = \qty{0.0100}{mol}$; $m = 0.0100 \times 159.6 =
\qty{1.60}{g}$.
\end{solution}

\begin{solution}{exo:g10:concentration-and-dilution:4}
$F = 1.0 / 0.050 = 20$; $V_0 = \qty{100.0}{mL} / 20 = \qty{5.0}{mL}$.
\end{solution}

\begin{solution}{exo:g10:concentration-and-dilution:5}
$C_m = c \times M = 0.10 \times 58.5 = \qty{5.85}{g/L}$, about
\qty{5.9}{g/L}.
\end{solution}

\begin{solution}{exo:g10:concentration-and-dilution:6}
$V_0 = c_1 V_1 / c_0 = 0.020 \times 250.0 / 0.50 = \qty{10.0}{mL}$
(\omterm{def:g10:concentration-and-dilution:dilution}{dilution factor} 25). A \qty{10.0}{mL} volumetric pipette and a
\qty{250.0}{mL} volumetric flask.
\end{solution}

\begin{solution}{exo:g10:concentration-and-dilution:7}
The volumetric flask has a single mark, made for one volume and
accurate to a small fraction of a percent; the graduations of a beaker
are only a rough guide. Likewise a volumetric pipette delivers one
volume very accurately, while a graduated cylinder is wider and read
less precisely.
\end{solution}

\begin{solution}{exo:g10:concentration-and-dilution:8}
Between \qty{0.02}{mol/L} and \qty{0.04}{mol/L}: about \qty{0.03}{mol/L}.
For more precision, add standards between these two, for example at
0.025, 0.030 and \qty{0.035}{mol/L}: a finer scale where it is needed.
\end{solution}

\begin{solution}{exo:g10:concentration-and-dilution:9}
$M(\ce{C6H12O6}) = \qty{180.0}{g/mol}$; $c = 50 / 180.0 =
\qty{0.28}{mol/L}$.
\end{solution}

\begin{solution}{exo:g10:concentration-and-dilution:10}
$M = \qty{342.0}{g/mol}$, so $n = \qty{1.00}{mol}$ and $c =
\qty{1.00}{mol/L}$. Diluted ten times: $c = \qty{0.100}{mol/L}$, that is
$C_m = 0.100 \times 342.0 = \qty{34.2}{g/L}$. A spoonful of
\qty{0.020}{L} holds $34.2 \times 0.020 = \qty{0.68}{g}$ of sucrose.
\end{solution}

\begin{solution}{exo:g10:concentration-and-dilution:11}
(a) The same mass of \omterm{def:g6:solutions-and-solubility:solute}{solute} in twice the volume: \qty{6.0}{g/L}.
(b) The volume is about halved, the mass of \omterm{def:g6:solutions-and-solubility:solute}{solute} unchanged: about
\qty{24}{g/L}.
\end{solution}

\begin{solution}{exo:g10:concentration-and-dilution:12}
\begin{enumerate}
  \item Three \omterm{def:g10:concentration-and-dilution:dilution}{dilutions} by 10 make a \omterm{def:g10:concentration-and-dilution:dilution}{dilution} by $10^3$:
        $\qty{0.10}{mol/L} / 1000 = \qty{1.0e-4}{mol/L}$.
  \item A single \omterm{def:g10:concentration-and-dilution:dilution}{dilution} by 1000 would need either a tiny volume of
        stock (\qty{0.1}{mL} into \qty{100}{mL}), impossible to measure
        accurately with a pipette, or a very large flask (\qty{1}{mL}
        into \qty{1}{L}). Three tenfold \omterm{def:g10:concentration-and-dilution:dilution}{dilutions} use ordinary pipettes
        and flasks.
\end{enumerate}
\end{solution}

\begin{solution}{exo:g10:concentration-and-dilution:13}
\begin{enumerate}
  \item $M(\ce{CaCl2}) = 40.1 + 2 \times 35.5 = \qty{111.1}{g/mol}$;
        $n = 11.1 / 111.1 = \qty{0.100}{mol}$; $c = 0.100 / 0.5000 =
        \qty{0.200}{mol/L}$.
  \item Each \ce{CaCl2} gives one \ce{Ca^2+} and two \ce{Cl-}:
        $[\ce{Ca^2+}] = \qty{0.200}{mol/L}$ and $[\ce{Cl-}] =
        \qty{0.400}{mol/L}$.
\end{enumerate}
\end{solution}

\begin{solution}{exo:g10:concentration-and-dilution:14}
\begin{enumerate}
  \item $r = \qty{0.60}{cm}$, $h = \qty{0.20}{cm}$:
        $V = \pi \times 0.60^2 \times 0.20 = \qty{0.23}{cm^3}$, that is
        \qty{0.23}{mL} too much water.
  \item The same \omterm{def:g6:solutions-and-solubility:solute}{solute} is in \qty{100.23}{mL} instead of \qty{100.00}{mL}:
        the concentration is too low by $0.23 / 100.23 \approx
        \qty{0.23}{\percent}$.
\end{enumerate}
\end{solution}

\begin{solution}{exo:g10:concentration-and-dilution:15}
$n = 0.20 \times 0.100 + 0.10 \times 0.300 = 0.020 + 0.030 =
\qty{0.050}{mol}$ in \qty{0.400}{L}: $c = \qty{0.125}{mol/L}$. It is not
the plain average, \qty{0.15}{mol/L}: there is three times more of the
dilute solution, so the \omterm{def:g6:pure-substances-and-mixtures:mixture}{mixture} is closer to \qty{0.10}{mol/L}. It is the
average weighted by the volumes.
\end{solution}

\begin{solution}{pb:g10:concentration-and-dilution:1}
\textbf{1.} \qty{0.9}{g} per \qty{0.100}{L}: $C_m = \qty{9.0}{g/L}$.

\textbf{2.} \qty{9.0}{g} per \qty{1000}{mL}, that is \qty{9.0}{mg} per
millilitre.

\textbf{3.} $m = 9.0 \times 0.500 = \qty{4.5}{g}$.

\textbf{4.} Yes: about \qty{9}{g} per litre against some \qty{360}{g} per
kilogram of water at saturation, about forty times less.

\textbf{5.} $9.0 \times 1.0 = \qty{9.0}{g}$ of salt.

\textbf{6.} $M(\ce{NaCl}) = 23.0 + 35.5 = \qty{58.5}{g/mol}$.

\textbf{7.} $c = 9.0 / 58.5 = \qty{0.154}{mol/L}$.

\textbf{8.} $n = 0.154 \times 0.500 = \qty{0.0769}{mol}$ (or
$4.5 / 58.5$).

\textbf{9.} Each \ce{NaCl} gives one \ce{Na+} and one \ce{Cl-}: both at
\qty{0.154}{mol/L}.

\textbf{10.} $0.0769 \times \num{6.02e23} = \num{4.63e22}$ sodium \omterm{def:g9:ions:ion}{ions}.

\textbf{11.} $F = 200 / 9.0 = 22.2$.

\textbf{12.} The \qty{4.5}{g} of the bag all come from the concentrated
solution: $V_0 = \qty{4.5}{g} / \qty{200}{g/L} = \qty{0.0225}{L} =
\qty{22.5}{mL}$.

\textbf{13.} $V_0 = V_1 / F = 500 / 22.2 = \qty{22.5}{mL}$.

\textbf{14.} $c_0 = 200 / 58.5 = \qty{3.42}{mol/L}$;
$c_0 V_0 = 3.42 \times 0.0225 = \qty{0.0769}{mol}$ and
$c_1 V_1 = 0.154 \times 0.500 = \qty{0.0769}{mol}$: equal.

\textbf{15.} About $500 - 22.5 = \qty{478}{mL}$ of sterile water.

\textbf{16.} A graduated pipette of \qty{25}{mL} (or a burette), read to
\qty{0.1}{mL}.

\textbf{17.} In a \qty{500.0}{mL} volumetric flask, made up to the mark.

\textbf{18.} $C_m = 200 \times 23.0 / 500 = \qty{9.20}{g/L}$, too
concentrated by $0.20 / 9.0 \approx \qty{2.2}{\percent}$.

\textbf{19.} \qty{22.5}{mL} of the concentrated (``20 \%'') solution for
one \qty{500}{mL} bag.
\end{solution}
